\section{Facts and Existence}

\begin{quotex}
Existence as reality is not something static; it perpetually involves itself in a dynamic movement of ontological self-expression, articulating its infinite inner possibilities in gradations from the less determinate to the more determinate until it appears at the level of concrete forms, such that the particular existences which we regard as multiple and diverse `things' having separate individual `essences' are nothing but the modes and aspects of the reality of existence.
\flright{\textsc{Syed Muhammead Naquib al-Attas}, \emph{Prolegomena to the Metaphysics of Islam}}
\end{quotex}


Readers should meditate on the beauty of the above passage before continuing. Existence as reality is distinguished from existence as concept. Existence is not static; it is always in the process of the self-expression of its possibilities of manifestation. Existence descends from its essential form through the degrees of existence into more concreate forms.

Truth is the conformity to reality; the former is the logical aspect and the latter, the ontological aspect. Reality is always in process, so that

\begin{quotex}
the realities comprising the external world are not immediately or directly given in consciousness, but are abstractions of them in varying degrees performed by the external and internal senses, and transformed into knowledge of the external world by means of intellectual construction.
\end{quotex}


The corollary is that facts by themselves are insufficient for true knowledge. Facts presume that reality is static and that ``facts speak for themselves'', ignoring the larger picture. \textbf{Alan Watts} gives the following example in \emph{The Book on The Taboo Against Knowing Who You Are}:

\begin{quotationx}
Here is someone who has never seen a cat. He is looking through a narrow slit in a fence, and, on the other side, a cat walks by. He sees first the head, then the less distinctly shaped furry trunk, and then the tail. Extraordinary! The cat turns round and walks back, and again he sees the head, and a little later the tail. This sequence begins to look like something regular and reliable. Yet again, the cat turns round, and he witnesses the same regular sequence: first the head, and later the tail. Thereupon he reasons that the event head is the invariable and necessary cause of the event tail, which is the head's effect. This absurd and confusing gobbledygook comes from his failure to see that head and tail go together: they are all one cat.

The cat wasn't born as a head which, sometime later, caused a tail; it was born all of a piece, a head-tailed cat. Our observer's trouble was that he was watching it through a narrow slit, and couldn't see the whole cat at once.
\end{quotationx}


Beware of fact checkers with single vision.

\paragraph{Postscript}


Lest you assume that I am relying on unreliable sources, you need to consider the whole reality, not just the fact. The brother of a friend used to hand out copies of Watts' book to whoever would take it. I was befuddled, never having read anything like that before. I went on to devour his books, especially the theological books written prior to his hippie days in Sausalito. That was my introduction to metaphysics and symbolism. Moreover, he introduced me to Rene Guenon, whose books I then sought out. Unlike today, they were hard to come by. Effort makes you value something more highly, whereas today Guenon is often a mere commodity, just another piece in the perennialist puzzle.

\emph{When you feel the breath of an angel on the back of your neck, seize the opportunity.}


\flrightit{Posted on 2021-09-22 by Cologero}

